% Event-timing-focused revised introduction for Overleaf.
% Citations are written as author-year placeholders for later integration.

\documentclass[11pt]{article}
\usepackage[margin=1in]{geometry}
\usepackage{amsmath}
\usepackage[version=4]{mhchem}
\usepackage[T1]{fontenc}
\usepackage{lmodern}

\title{Precession-Phase Organization of DO-Like Asian Monsoon Transitions}
\author{}
\date{}

\begin{document}
\maketitle
\maketitle
\section*{Introduction}

Millennial-scale climate variability (MCV) is a major expression of abrupt climate change in glacial climate systems. Its canonical manifestation is the Dansgaard--Oeschger (DO) sequence in Greenland ice cores, where rapid stadial--interstadial transitions record major reorganizations of the North Atlantic climate system \citep{rasmussen2014}. Yet MCV was not confined to Greenland. Marine, ice-core, and speleothem archives show that abrupt millennial variability occurred across multiple Pleistocene glacial cycles and was expressed in both high-latitude and low-latitude climate systems \citep{barker2011, sun2021, hodell2023}.

Speleothem records provide a particularly powerful archive because their U--Th chronologies allow abrupt hydroclimate changes to be compared directly with ice-core event sequences. During the last glacial period, independently dated speleothem records indicate that abrupt Greenland warmings were accompanied by near-synchronous hydroclimate shifts in monsoon and mid-latitude regions \citep{corrick2020}. These shifts in speleothem record can therefore be treated as precisely dated DO-like Asian monsoon expressions of MCV.

The intensity of MCV is strongly conditioned by the background climate state, with higher MCV activity during intermediate glacial state \citep{dansgaard1993, ganopolski2001, kawamura2017, zhang2017}. Long proxy records further suggest that the magnitude of MCV was modulated by orbital forcing, with precession and obliquity leaving detectable signatures in the millennial-scale variance envelope \citep{sun2021, hodell2023, thirumalai2020}. 

Amplitude modulation, however, does not necessarily imply event pacing. A larger millennial-scale envelope under a particular orbital configuration may reflect the amplification of internally generated transitions, rather than external control on when individual events occur. In stochastic or excitable climate systems, abrupt regime shifts can arise when internal variability crosses a threshold, producing clustered or intermittent events without direct orbital pacing \citep{ditlevsen2005, kleppin2015, lohmann2018}. Conversely, climate-model experiments suggest that orbital configuration can alter DO-like oscillatory regimes by changing tropical hydroclimate, high-latitude sea ice, and ocean circulation, thereby switching millennial-scale AMOC variability on or off under otherwise comparable boundary conditions \citep{zhang2021, kuniyoshi2022, malmierca2023}. The key unresolved proxy question is therefore not only whether orbital forcing modulates MCV amplitude, but whether it also organizes the timing of individual stadial--interstadial-like transitions.

Testing this event-timing question requires a long, precisely dated sequence of abrupt transitions. The 640 kyr Chinese speleothem composite of \citet{cheng2016}, together with the objectively identified weak- and strong-monsoon starts of \citet{rousseau2023}, provides a rare opportunity to examine event timing over multiple orbital cycles. These weak and strong starts are not Greenland-defined DO events themselves, but DO-like Asian monsoon transitions that record the monsoon-sector expression of MCV.

Here we test whether the timing of these DO-like Asian monsoon transitions is organized by orbital phase. Using the 0--640 ka weak- and strong-monsoon start catalogues of \citet{rousseau2023}, we ask three linked questions. First, are the event sequences consistent with stationary random occurrence? Second, do weak- and strong-monsoon starts cluster at preferred precession or obliquity phases? Third, does any orbital-phase dependence persist after accounting for slow background states represented by LR04 ice volume and atmospheric \(\ce{CO2}\)? By shifting the focus from variance amplitude to transition probability, this event-based framework tests whether orbital configuration not only modulated the strength of MCV, but also helped organize the timing of abrupt monsoon-state changes.

\begin{figure}[t]
    \centering
    \includegraphics[width=\textwidth]{figures/fig01.pdf}
    \caption{
    \textbf{Data and phase framework.}
    (a) Chinese speleothem $\delta^{18}\mathrm{O}$ composite with strong and weak monsoon-start timings marked by upper ticks.
    (b) LR04 benthic $\delta^{18}\mathrm{O}$.
    (c) Atmospheric CO$_2$.
    (d) Precession phase used in the hazard model, shown with the raw precession index.
    }
    \label{fig:data_phase_framework}
\end{figure}



\section*{Methods}

We analyzed the weak and strong monsoon-start chronologies identified by Rousseau et al. (2023) from the Chinese speleothem composite of Cheng et al. (2016). There are 96 strong monsoon starts and 94 weak monsoon starts. All tests below were applied independently to the two event catalogues.

Before fitting the main event-hazard models, we performed two supporting event-timing tests. First, we used a Lohmann-style stationary occurrence test to ask whether the weak and strong start sequences are compatible with stationary random timing. Second, we used circular Rayleigh tests to ask whether event phases are uniformly distributed with respect to precession and obliquity. These tests provide diagnostic context for the hazard analysis, but they do not adjust for glacial background state; full details are given in Text S1 and Text S2 of the Supporting Information.

The main analysis used a binned Poisson hazard generalized linear model to test whether orbital phase explains transition probability after accounting for glacial background state. The 0--640 ka interval was divided into bins of width \(\Delta t = 0.2\) kyr. For each event catalogue and bin \(i\), the response variable was the event count \(Y_i\). Conditional on the covariates, event counts were modeled as independent Poisson variables:
\[
Y_i \sim \mathrm{Poisson}(\mu_i),
\qquad
\mu_i = \Delta t_i \lambda_i,
\]
where \(\lambda_i\) is the fitted event rate per kyr. The log link was applied to the rate rather than directly to the expected bin count:
\[
\log \lambda_i
= \beta_0 + \mathbf{x}_i^\mathsf{T}\boldsymbol{\beta}.
\]
Equivalently,
\[
\log \mu_i
= \log \Delta t_i + \beta_0 + \mathbf{x}_i^\mathsf{T}\boldsymbol{\beta},
\]
so \(\log\Delta t_i\) enters the model as an offset. The Poisson log-likelihood was
\[
\ell(\boldsymbol{\beta})
=
\sum_i
\left[
Y_i\log\mu_i - \mu_i - \log(Y_i!)
\right].
\]

The baseline model set consisted of four candidate models: a stationary intercept-only model, a climate-background model with LR04 and atmospheric \(\ce{CO2}\), a precession-phase-only model, and a combined model:
\[
\log \lambda_i
= \beta_0
+ \beta_1 \mathrm{LR04}_i
+ \beta_2 \ce{CO2}_i
+ \beta_3 \sin\theta_i
+ \beta_4 \cos\theta_i,
\]
where \(\theta_i\) is the precession phase at the bin center. Precession phase was defined from the precession-index waveform by assigning phase \(0\) to local minima and phase \(\pi\) to local maxima, with linear interpolation of unwrapped phase between adjacent extrema. LR04 and \(\ce{CO2}\) were interpolated to bin centers and range-scaled to have zero mean and unit observed range. Precession phase was represented by \(\sin\theta_i\) and \(\cos\theta_i\), because phase is circular and cannot be treated as a linear predictor.

Models with predictors were fitted by maximum likelihood using L-BFGS-B optimization; the stationary model has the analytic maximum-likelihood rate \(N/T\). Models were ranked using sample-size-corrected Akaike information criterion (AICc). Nested models were compared with likelihood-ratio tests, using
\[
\mathrm{LR} = 2(\ell_\mathrm{full}-\ell_\mathrm{reduced})
\]
and a chi-square reference distribution with degrees of freedom equal to the number of added coefficients. The key comparison was ``precession phase after climate'', in which the reduced model was LR04 + \(\ce{CO2}\) and the full model added the two circular precession-phase terms.

For models containing \(\sin\theta\) and \(\cos\theta\), we converted the two phase coefficients into an interpretable preferred phase. If the phase component is
\[
\beta_s\sin\theta+\beta_c\cos\theta,
\]
then the phase of maximum fitted rate is
\[
\theta^\ast = \operatorname{atan2}(\beta_s,\beta_c)\ \mathrm{mod}\ 2\pi,
\]
and the phase-response amplitude is
\[
A = (\beta_s^2+\beta_c^2)^{1/2}.
\]
The fitted maximum-to-minimum rate ratio over one precession cycle is therefore \(\exp(2A)\). 

We performed two sensitivity experiments for the hazard-model results. First, we asked whether other plausible forcings add explanatory power beyond the baseline LR04 + \(\ce{CO2}\) + precession-phase model. The extra predictors were Antarctic temperature (AT), obliquity, eccentricity, and \(65^\circ\)N summer-solstice insolation. Each extra predictor was range-scaled in the same way as LR04 and \(\ce{CO2}\). We tested the contribution of the extra forcings in three ways: adding each extra predictor to the baseline model one at a time, adding all four extra predictors jointly, and dropping each extra predictor from the full extended model to estimate its unique contribution after the other extra predictors were present.

Second, we tested sensitivity to the bin width used in the binned hazard approximation. The baseline bin width was 0.2 kyr, chosen to keep bins short relative to expected changes in the underlying hazard. To check that this choice did not determine the conclusions, we repeated the core model comparison for bin widths of 0.2, 0.4, 0.6, 0.8, and 1.0 kyr. For each bin width, we tracked model ranking by AICc, the likelihood-ratio \(p\) value for precession phase after LR04 + \(\ce{CO2}\), the preferred precession phase, and the sparsity of the event-count bins.

\begin{figure}[t]
    \centering
    \includegraphics[width=\textwidth]{figures/fig02.pdf}
    \caption{
    \textbf{Event-timing randomness and precession-phase clustering.}
    (a,b) Lohmann-style stationary-event tests for strong and weak monsoon starts, comparing the observed $E_S$ statistic with stationary Poisson and fixed-$N$ null distributions.
    (c,d) Rayleigh tests for clustering of strong and weak monsoon starts in precession phase; arrows show the mean resultant direction and dashed circles mark the $p=0.05$ threshold.
    }
    \label{fig:randomness_rayleigh}
\end{figure}



\section*{Results}

The supporting event-timing tests motivate a non-stationary, phase-aware hazard model. The Lohmann-style test shows that the weak monsoon-start sequence is inconsistent with stationary random occurrence under both the unconditional Poisson and fixed-\(N\) nulls, whereas the strong monsoon-start sequence is significant under the fixed-\(N\) null and marginal under the unconditional Poisson null (Figure 2). Rayleigh tests further show significant precession-phase clustering for both strong and weak starts, with perfered phase for the strong and weak monsoon starts to be 21.9° and 226.7°, respectively (Figure 2). Conversely, no significant obliquity-phase clustering. Detailed statistics and test definitions are reported in Text S1 and Text S2 of the Supporting Information.

The binned Poisson hazard models confirm that precession phase contributes information beyond glacial background state. For both weak and strong monsoon starts, the combined LR04 + \(\ce{CO2}\) + precession-phase model had the lowest AICc among the core candidate models. For strong monsoon starts, adding LR04 + \(\ce{CO2}\) to the stationary model improved the likelihood (\(p=0.0092\)), and adding precession phase to LR04 + \(\ce{CO2}\) improved it further (\(\mathrm{LR}=15.10\), \(p=5.25\times10^{-4}\), \(\Delta\mathrm{AICc}=-11.09\)). The fitted preferred precession phase in the full model was \(10.8^\circ\), close to the Rayleigh mean direction, and the fitted maximum-to-minimum rate ratio over the precession cycle was 3.18, which implies an approximately threefold contrast in event hazard across the precession cycle.

Weak monsoon starts show the same model-ranking pattern but with the opposite preferred precession sector. The LR04 + \(\ce{CO2}\) model improved on stationarity (\(p=4.31\times10^{-4}\)), and precession phase still improved the model after LR04 + \(\ce{CO2}\) were included (\(\mathrm{LR}=10.66\), \(p=0.00485\), \(\Delta\mathrm{AICc}=-6.65\)). The full model preferred a precession phase of \(237.5^\circ\), again close to the Rayleigh mean direction for weak starts, with a fitted maximum-to-minimum rate ratio of 2.63. These results indicate that precession phase is not merely reproducing the slow LR04/\(\ce{CO2}\) background structure; it remains associated with event probability after those background variables are included. 


The extra-forcing sensitivity experiment does not identify an additional independent driver comparable to precession phase (Figure S*, S*). After the baseline LR04 + \(\ce{CO2}\) + precession-phase model was fitted, adding AT, obliquity, eccentricity, or \(65^\circ\)N summer-solstice insolation one at a time did not significantly improve either event model. For strong monsoon starts, the single-extra \(p\) values ranged from 0.227 to 0.936; for weak monsoon starts, they ranged from 0.363 to 0.689. Adding all four extra forcings jointly also failed to improve the baseline model for strong starts (\(p=0.682\), \(\Delta\mathrm{AICc}=+5.75\)) or weak starts (\(p=0.733\), \(\Delta\mathrm{AICc}=+6.02\)). Drop-one tests within the full extended model likewise showed no significant unique contribution from any extra forcing. Thus, in this event-hazard framework, the robust orbital signal is the precession phase term already present in the baseline model, rather than obliquity, eccentricity, Antarctic temperature, or \(65^\circ\)N summer-solstice insolation as independent additions.

The bin-width sensitivity experiment shows that the 0.2 kyr bin choice is not driving the qualitative conclusions (Table S*). Across bin widths from 0.2 to 1.0 kyr, the LR04 + \(\ce{CO2}\) + precession-phase model remained the best core model by AICc for both event catalogues. The likelihood-ratio test for precession phase after LR04 + \(\ce{CO2}\) stayed significant at every tested bin width. For strong monsoon starts, the preferred phase stayed between \(10.5^\circ\) and \(11.5^\circ\), with rate ratios of approximately 3.15--3.19. For weak monsoon starts, the preferred phase stayed between \(236.6^\circ\) and \(238.4^\circ\), with rate ratios of approximately 2.63--2.69. The preferred phases shifted by less than \(1^\circ\) relative to the 0.2 kyr baseline. The 0.2 kyr bins are sparse, with about 97\% empty bins and no multi-event bins, but wider bins also preserve the same phase preference and model ranking. We therefore treat 0.2 kyr as a conservative baseline that keeps the within-bin constant-hazard approximation short while yielding conclusions that are stable across coarser binning choices.

\section*{Discussion}

Our results suggest that the abrupt monsoon transitions identified in the Chinese speleothem composite are not simply embedded within an orbitally modulated background record; their event timings themselves are organized by precession phase. This distinction is important because precession pacing of continuous Asian speleothem \(\delta^{18}\mathrm{O}\) variability is already well established \citep{wang2008, cheng2016}. The question addressed here is narrower and more event-based: given a catalogue of objectively identified transitions toward strong and weak monsoon states, is the probability of transition elevated during particular portions of the precession cycle, and does this relationship remain after accounting for slow background state? The agreement between the non-parametric Rayleigh tests and the binned Poisson hazard models supports such an interpretation. Strong monsoon starts cluster near the precession-index minimum sector, whereas weak monsoon starts cluster closer to the opposite sector. The fitted hazard model recovers nearly the same preferred phases after LR04 and \(\ce{CO2}\) are included, indicating that the phase signal is not only a consequence of the long glacial background trend.

A useful conceptual interpretation is that glacial background state controls the susceptibility of the climate system, whereas precession phase provides a seasonal-orbital window in which transitions are more likely. LR04 and \(\ce{CO2}\) summarize slow changes in ice volume, greenhouse forcing, sea level, boundary conditions, and ocean-atmosphere background state. These variables plausibly affect whether the coupled monsoon--North Atlantic system is close to a threshold or in a regime capable of millennial-scale oscillations. Precession phase, in contrast, describes where the system lies within the low-latitude seasonal-insolation cycle. In the fitted model, both components matter: LR04 + \(\ce{CO2}\) improves on the stationary model, but adding precession phase further improves the likelihood and lowers AICc for both event types. This favors a ``background susceptibility plus orbital phase organization'' view rather than either a purely stochastic event sequence or a purely background-state-controlled sequence.

This interpretation is consistent with recent moisture-source modeling for the East Asian monsoon. \citet{xie2023} used transient water-tagging simulations to show that northern East Asian monsoon-region precipitation has a strong 23 kyr component primarily controlled by precession-induced Northern Hemisphere summer insolation, while global ice volume modulates precipitation through changes in terrestrial and Pacific moisture-source contributions. Their mechanism is not an event-timing model, but it maps closely onto the statistical structure found here: precession phase is the dominant orbital term, while ice-volume and greenhouse-background variables cannot be ignored. In this sense, our GLM does not imply that precession acts in isolation. Rather, it suggests that precession-phase-dependent monsoon-season hydroclimate and moisture transport may define favorable transition windows, with LR04/\(\ce{CO2}\)-related background state modulating how often those windows produce abrupt transitions.

The opposite preferred phases for strong and weak monsoon starts also carry physical information. Strong starts, which correspond to transitions toward stronger Asian monsoon states and Greenland interstadial-like conditions, peak near \(10^\circ\)--\(22^\circ\) in our phase convention. Because precession-index minima are assigned to \(0^\circ\), this phase is close to the sector associated with stronger Northern Hemisphere summer insolation. Weak starts peak near \(227^\circ\)--\(238^\circ\), approximately the opposite precession sector. This antiphased behavior is unlikely to arise if the event catalogue were merely sampling a single broad interval of enhanced millennial variability. Instead, it suggests that transitions into and out of strong monsoon states are asymmetric with respect to the precession cycle. In a threshold framework, the same slowly varying background state may permit abrupt transitions, but the direction of transition may depend on whether precession favors monsoon strengthening or weakening.

The absence of an independent obliquity signal in the event-timing tests should be interpreted cautiously. Obliquity has been implicated in the amplitude modulation of millennial-scale variability and in model simulations of DO-like oscillations \citep{sun2021, hodell2023, zhang2021, malmierca2023}. Our result is more specific: in this event catalogue and within the present hazard-model design, obliquity phase does not show significant circular clustering, and adding obliquity as an additional predictor does not significantly improve the LR04 + \(\ce{CO2}\) + precession-phase baseline. This does not rule out obliquity as a background control on ice sheets, sea ice, or MCV amplitude. It means that the timing of these speleothem-recorded transitions contains a more clearly separable precession-phase signal than an obliquity-phase signal.

Similarly, the lack of additional explanatory power from \(65^\circ\)N summer-solstice insolation should not be read as evidence that summer insolation is unimportant. A scalar insolation time series and a circular precession-phase representation ask different statistical questions. The \(65^\circ\)N summer-solstice series combines precession, obliquity, and eccentricity effects at a fixed latitude and calendar day, whereas the circular phase terms test whether event probability depends on position within the precession cycle, regardless of the absolute amplitude of the insolation curve. Recent model work emphasizes that precession affects not only peak summer insolation but also seasonal transitions, monsoon onset, retreat, and moisture transport pathways \citep{wu2023}. Thus, the stronger performance of precession phase may indicate that the relevant event-timing information is phase-like and seasonal, rather than captured by a single high-latitude summer-solstice forcing index.

The results also clarify how this study relates to broader MCV and DO-event literature. The events analyzed here are not Greenland-defined DO events for the full 0--640 ka interval. They are Chinese speleothem-recorded abrupt monsoon transitions, objectively identified from the Cheng et al. composite by \citet{rousseau2023}. For the last glacial period, speleothem abrupt changes are tightly linked to Greenland stadial--interstadial transitions \citep{corrick2020}, supporting their use as a monsoon-sector expression of MCV. Beyond the last glacial period, however, the terminology must remain cautious: these are DO-like Asian monsoon transitions, not directly observed Greenland DO events. This distinction is not a weakness of the dataset; it is what makes the record useful. The speleothem composite provides a long, U--Th-dated event sequence through multiple orbital cycles, allowing event-timing questions that cannot be addressed using the last glacial Greenland sequence alone.

The lagged predictive-information analysis provides an additional, exploratory timescale perspective. Because the event series is sparse, we did not use a fully non-parametric transfer-entropy estimator. Instead, we used nested Poisson hazard models to ask how much a lagged driver improves prediction of binned event occurrence after conditioning on recent event history. For a driver \(X\) and lag \(\tau\), the reduced model included the 2 kyr event-history term, and the full model added \(X(t+\tau)\), where age increases into the past. The log-likelihood gain,
\[
\Delta \ell(\tau)
= \ell(Y_t \mid Y_{\mathrm{hist}}, X_{t+\tau})
- \ell(Y_t \mid Y_{\mathrm{hist}}),
\]
was converted to bits per bin and bits per event, with likelihood-ratio \(p\) values and Benjamini--Hochberg \(q\) values used as diagnostic statistics. This quantity should be interpreted as model-based lagged predictive information, not as proof of strict causal transfer entropy.

The lagged information results are nevertheless consistent with the hazard-model interpretation. For strong monsoon starts, the best single-driver lags for LR04 and \(\ce{CO2}\) are both about 4 kyr, whereas precession phase peaks at 0 kyr. For weak monsoon starts, the best lags are shorter, with \(\ce{CO2}\) peaking near 0.6 kyr and LR04 and precession phase near 0 kyr. In the stricter conditional tests, precession phase still adds information after LR04 and \(\ce{CO2}\) are included at the same lag or at their best single-driver lags. The strongest conditional precession-phase signal occurs near zero lag for strong starts and near zero to sub-kyr lag for weak starts. This pattern is compatible with precession phase acting as a transition-window variable, while LR04 and \(\ce{CO2}\) represent slower background susceptibility. Because the analysis is based on binned event counts and interpolated forcings, however, these lag estimates should be viewed as order-of-magnitude diagnostics rather than precise response times.

Several limitations should be kept in view. First, the event catalogue depends on the Rousseau et al. transition-detection procedure applied to a composite speleothem record. Different detection thresholds, filtering choices, or age-model ensembles could slightly alter individual event timings. The stability of the preferred phase across bin widths is encouraging, but it does not replace a full age-uncertainty propagation. Second, speleothem \(\delta^{18}\mathrm{O}\) is not a direct local precipitation amount. It integrates upstream rainout, moisture-source changes, circulation pathways, seasonality, and regional hydroclimate effects \citep{liu2014, cheng2021, wang2022}. The terms ``strong'' and ``weak'' monsoon states are therefore best understood as large-scale speleothem-recorded monsoon-system states rather than simple local rainfall anomalies. This caveat is particularly relevant in light of moisture-source results such as \citet{xie2023}, which show that orbital precipitation changes can involve compensating terrestrial and oceanic source contributions.

Third, the Poisson hazard model is an intentionally simple statistical approximation. The 0.2 kyr bin width keeps the within-bin constant-hazard assumption conservative, and the bin-width sensitivity test shows that the qualitative results are stable up to 1 kyr. Still, the model assumes conditional independence of bins given covariates, does not explicitly represent event-duration dynamics, and uses a limited set of background predictors. The likelihood-ratio tests therefore establish statistical association and predictive improvement, not a complete dynamical mechanism. Fourth, covariates such as precession phase, summer insolation, eccentricity, and ice-volume-related background state are not independent components of the climate system. The sensitivity results show that additional variables do not improve the selected baseline in this model, but they should not be interpreted as excluding indirect pathways through which obliquity, eccentricity, or ice sheets influence MCV.

Despite these caveats, the event-based signal is robust across several complementary tests. The stationary tests show that the transition sequences are not fully compatible with simple random occurrence. The Rayleigh tests show direct precession-phase clustering. The Poisson hazard GLM shows that precession phase remains informative after LR04 and \(\ce{CO2}\) are included. The sensitivity experiments show that this conclusion is stable to bin width and is not replaced by AT, obliquity, eccentricity, or \(65^\circ\)N summer-solstice insolation as independent additions. Together, these results point to precession-phase organization of abrupt monsoon-transition probability, superimposed on slower glacial background modulation.

\section*{Conclusions}

This study tests whether the timing of DO-like Asian monsoon transitions in the 0--640 ka Chinese speleothem composite is organized by orbital phase. The main result is that strong and weak monsoon starts show significant and opposite precession-phase preferences, and that these preferences persist in a binned Poisson hazard model after accounting for LR04 ice-volume background and atmospheric \(\ce{CO2}\). Strong starts preferentially occur near the precession-index minimum sector, while weak starts occur closer to the opposite precession sector. The combined LR04 + \(\ce{CO2}\) + precession-phase model is the best core model for both event types, and the inferred phase preferences are stable across bin-width sensitivity tests. Additional forcings, including obliquity, eccentricity, Antarctic temperature, and \(65^\circ\)N summer-solstice insolation, do not add comparable independent predictive power in this framework. A complementary lagged predictive-information analysis suggests that the precession-phase contribution is strongest at near-zero lag, whereas LR04 and \(\ce{CO2}\) behave more like background-state controls. We therefore interpret the results as evidence that precession phase helps organize the timing of speleothem-recorded abrupt monsoon transitions, while glacial background state modulates the system's susceptibility to such transitions. This event-timing perspective extends earlier work on orbital modulation of continuous speleothem \(\delta^{18}\mathrm{O}\) and MCV amplitude, and provides a statistical target for transient climate-model experiments that seek to explain how orbital forcing and glacial boundary conditions interact to pace abrupt hydroclimate changes.

\end{document}


